\begin{abstract}
Large language models (LLMs) can translate SAP Adaptive Server Enterprise (ASE) Transact-SQL to
Microsoft SQL Server, but a translation that compiles may still behave differently. We present VeriMig,
a framework that checks each translation with four gates: deployment, a lint for remaining ASE
constructs, differential execution against a reviewed reference translation, and data reconciliation.
Only the first two gates give feedback to the model's repair loop, so the references never shape the
translations they score. We contribute a taxonomy of \NumCategories{} categories of dialect differences,
a corpus of \NumDevCases{} development cases and \NumHeldout{} held-out procedures, and an evaluation of
\NumModels{} locally hosted code models. Without an ASE server, expected behaviour comes from reference
translations; \NumReviewCases{} rely on documented, not executed, ASE semantics.
For the strongest local model, agreement with the references rises from \PlainEquiv\% with a plain prompt
to \TrapEquiv\% with a checklist of known dialect differences and \AgenticEquiv\% with the repair loop.
The checklist improves all five local models; the further gain from repair is not significant. Of the
plain-prompt translations, \PlainSilent\% compile and pass the lint yet disagree with the reference, and on
the held-out procedures agreement is \HeldoutPlainEquiv\% with the plain prompt and
\HeldoutAgenticEquiv\% with the repair loop. A hosted reference model reaches
\HostedPlainEquiv\% with the plain prompt and \HostedTrapEquiv\% with the checklist, so these failures are
mainly a property of smaller local models. A fault-injection study shows that only a per-row hash over an
unambiguous row encoding detects every tested data fault. All artifacts are released.
\end{abstract}

\begin{IEEEkeywords}
database migration, SQL dialect translation, large language models, differential testing,
Sybase ASE, SQL Server, verification, agentic AI
\end{IEEEkeywords}
